\seccionanexo{Anexo A: Instrumentos de Recolección de Datos}

Cuestionario estructurado en escala tipo Likert de 5 puntos (1: Totalmente en desacuerdo a 5: Totalmente de acuerdo) aplicado a la muestra de 60 estudiantes postulantes en el MÓDULO TECNOLÓGICO PRODUCTIVO SAN JULIÁN (anonimizado). Tiempo estimado: 4 min.

\begin{enumerate}
    \item \textbf{Usabilidad y Diseño Móvil (Vista):}
    \begin{itemize}
        \item La interfaz web es clara y se adapta adecuadamente a la pantalla de mi teléfono celular.
        \item Los pasos para completar el formulario de preinscripción son comprensibles y rápidos de ejecutar.
    \end{itemize}
    \item \textbf{Autoservicio y Recuperación de Credenciales (Controlador/Modelo):}
    \begin{itemize}
        \item Pude consultar mi usuario y restablecer mi contraseña de forma autónoma con mi cédula de identidad y fecha de nacimiento ($<$2 min).
        \item El mecanismo de autoservicio me permitió asegurar oportunamente mi cupo en la especialidad técnica de mi preferencia sin hacer filas presenciales.
    \end{itemize}
    \item \textbf{Matriculación Documental y Transparencia:}
    \begin{itemize}
        \item El comprobante digital emitido refleja con claridad los requisitos documentales exigidos (libretas de 1ro a 3ro, cédulas, carnet y certificado).
        \item La entrega y verificación de mis documentos con los tutores de curso se registró de manera ágil y ordenada en el sistema.
    \end{itemize}
    \item \textbf{Organización en Cursos y Paralelos:}
    \begin{itemize}
        \item La asignación a mi paralelo (4to A o 4to B) respetó la compatibilidad de horarios con mi colegio de origen.
        \item La distribución de plazas entre los colegios del distrito municipal fue equitativa y transparente.
    \end{itemize}
\end{enumerate}

\vspace{0.4cm}
\noindent\textbf{Guía de Entrevista Semiestructurada — Personal Administrativo (5 funcionarios: Director, Secretaría Académica y 3 docentes entre ellos Tutores de Curso):}

\begin{enumerate}
    \item ¿Qué dificultades enfrentaron en la gestión 2025 con los reseteos manuales de contraseñas y la pérdida de cupos vocacionales?
    \item ¿Cómo se recepcionaban y archivaban las libretas de 1ro a 3ro de secundaria, cédulas de tutores, carnet y certificado antes de la V2.0?
    \item ¿Qué porcentaje de expedientes quedaba incompleto o sin trazabilidad en 2025?
    \item ¿Qué errores y sobrecarga generaba la asignación manual a paralelos en planillas externas por unidad educativa?
    \item ¿Qué conflictos de coordinación horaria surgían entre el MTP y las unidades educativas de origen?
    \item ¿Qué requisitos de validación consideran imprescindibles para el módulo de matriculación documental auditada?
    \item ¿Qué criterios de equilibrio por taller deberían aplicarse en la asignación automatizada (4to A/B)?
    \item ¿Qué expectativas tienen sobre el panel de control y reportes en tiempo real para la toma de decisiones?
\end{enumerate}

\notatabla{Instrumentos aplicados en el primer bimestre de 2026; datos anonimizados (C.I. como ****). Validación por juicio del tutor y piloto con 5 estudiantes (ver Anexo B).}
